\documentclass[12pt, a4paper]{ctexart}

\usepackage{amsmath, amssymb, amsthm}
\usepackage{geometry}
\usepackage{fancyhdr}
\usepackage{enumitem}
\usepackage{xcolor}
\usepackage{hyperref}

\geometry{margin=1in}
\pagestyle{fancy}
\fancyhf{}
\rhead{线性代数与微分方程}
\lhead{Wang ChunYan}
\rfoot{Page \thepage}

\newcommand{\problem}[1]{\subsection*{第 #1 题}}
\newcommand{\solution}{\paragraph{解.}}

\definecolor{solutioncolor}{RGB}{0, 100, 0}

\begin{document}

\begin{center}
    {\LARGE\bfseries 线性代数与微分方程 作业解答} \\[0.5em]
    {\large 线性代数与微分方程} \\[0.3em]
    Wang ChunYan \quad | \quad 2026-10-04
\end{center}

\hrule
\vspace{1em}
\problem{1}

计算行列式

设三阶方阵

$$
A = \begin{pmatrix} 3 & 1 & 1 \\ 1 & 3 & 1 \\ 1 & 1 & 3 \end{pmatrix}
$$

求行列式 $\det(A)$ 的值。

---

\solution

对矩阵 $A = \begin{bmatrix}3 & 1 & 1 \\ 1 & 3 & 1 \\ 1 & 1 & 3\end{bmatrix}$ 计算行列式

$\det(A) = 20$

\textbf{Answer:}
\[\boxed{20}\]

\vspace{1em}
\hrule
\vspace{1em}

\problem{2}

求逆矩阵

设矩阵

$$
A = \begin{pmatrix} 2 & 1 \\ 1 & 1 \end{pmatrix}
$$

判断 $A$ 是否可逆；若可逆，求其逆矩阵 $A^{-1}$。

---

\solution

先判断可逆性：$\det(A) = 1$

$\det(A) \neq 0$，故可逆。

$A^{-1} = \begin{bmatrix}1 & -1 \\ -1 & 2\end{bmatrix}$

\textbf{Answer:}
\[\begin{bmatrix}1 & -1 \\ -1 & 2\end{bmatrix}\]

\vspace{1em}
\hrule
\vspace{1em}

\problem{3}

求解线性方程组

求解线性方程组 $Ax = b$，其中

$$
A = \begin{pmatrix} 1 & 1 & 1 \\ 2 & -1 & 1 \\ 1 & 2 & -1 \end{pmatrix}, \quad
b = \begin{pmatrix} 6 \\ 3 \\ 2 \end{pmatrix}
$$

即：

$$
\begin{cases} x + y + z = 6 \\ 2x - y + z = 3 \\ x + 2y - z = 2 \end{cases}
$$

---

\solution

系数矩阵 $A = \begin{bmatrix}1 & 1 & 1 \\ 2 & -1 & 1 \\ 1 & 2 & -1\end{bmatrix}$，右端 $b = \begin{bmatrix}6 \\ 3 \\ 2\end{bmatrix}$

增广矩阵行最简形：$\left[\begin{matrix}1 & 0 & 0 & 1\\0 & 1 & 0 & 2\\0 & 0 & 1 & 3\end{matrix}\right]$

唯一解：$x = 1,\ y = 2,\ z = 3$

\textbf{Answer:}
\[\boxed{x = 1,\ y = 2,\ z = 3}\]

\vspace{1em}
\hrule
\vspace{1em}

\problem{4}

求特征值与特征向量

设矩阵

$$
A = \begin{pmatrix} 2 & 1 \\ 1 & 2 \end{pmatrix}
$$

求 $A$ 的全部特征值与对应的特征向量。

---

\solution

特征多项式：$\det(A - \lambda I) = \lambda^{2} - 4 \lambda + 3 = 0$

特征值 $\lambda = 1$（代数重数 1）

对应特征向量：$\begin{bmatrix}-1 \\ 1\end{bmatrix}$

特征值 $\lambda = 3$（代数重数 1）

对应特征向量：$\begin{bmatrix}1 \\ 1\end{bmatrix}$

\textbf{Answer:}
\[\boxed{\lambda = 1,\ \lambda = 3}\]

\vspace{1em}
\hrule
\vspace{1em}

\problem{5}

矩阵乘法

设

$$
A = \begin{pmatrix} 1 & 2 \\ 3 & 4 \end{pmatrix}, \quad
B = \begin{pmatrix} 5 & 6 \\ 7 & 8 \end{pmatrix}
$$

计算矩阵乘积 $AB$。

---

\solution

$A = \begin{bmatrix}1 & 2 \\ 3 & 4\end{bmatrix}$，$B = \begin{bmatrix}5 & 6 \\ 7 & 8\end{bmatrix}$

$AB = \begin{bmatrix}19 & 22 \\ 43 & 50\end{bmatrix}$

\textbf{Answer:}
\[\begin{bmatrix}19 & 22 \\ 43 & 50\end{bmatrix}\]

\vspace{1em}
\hrule
\vspace{1em}

\problem{6}

求矩阵的秩

设矩阵

$$
A = \begin{pmatrix} 1 & 2 & 3 \\ 4 & 5 & 6 \\ 7 & 8 & 9 \end{pmatrix}
$$

求 $\operatorname{rank}(A)$。

---

\solution

对矩阵 $A = \begin{bmatrix}1 & 2 & 3 \\ 4 & 5 & 6 \\ 7 & 8 & 9\end{bmatrix}$ 做初等行变换

行阶梯形：$\left[\begin{matrix}1 & 0 & -1\\0 & 1 & 2\\0 & 0 & 0\end{matrix}\right]$

$\operatorname{rank}(A) = 2$

\textbf{Answer:}
\[\boxed{2}\]

\vspace{1em}
\hrule
\vspace{1em}

\problem{7}

一阶线性微分方程（初值问题）

求解初值问题：

$$
y' + y = e^{x}, \qquad y(0) = 1
$$

---

\solution

微分方程：$y{\left(x \right)} - e^{x} + \frac{d}{d x} y{\left(x \right)} = 0$

初值条件：$y{\left(0 \right)} = 1$

解：$y(x) = \cosh{\left(x \right)}$

\textbf{Answer:}
\[\boxed{y(x) = \cosh{\left(x \right)}}\]

\vspace{1em}
\hrule
\vspace{1em}

\problem{8}

二阶常系数齐次微分方程（初值问题）

求解初值问题：

$$
y'' - 3y' + 2y = 0, \qquad y(0) = 1, \quad y'(0) = 0
$$

---

\solution

微分方程：$2 y{\left(x \right)} - 3 \frac{d}{d x} y{\left(x \right)} + \frac{d^{2}}{d x^{2}} y{\left(x \right)} = 0$

初值条件：$\left. \frac{d}{d x} y{\left(x \right)} \right|_{\substack{ x=0 }} = 0$

初值条件：$y{\left(0 \right)} = 1$

解：$y(x) = \left(2 - e^{x}\right) e^{x}$

\textbf{Answer:}
\[\boxed{y(x) = \left(2 - e^{x}\right) e^{x}}\]

\vspace{1em}
\hrule
\vspace{1em}

\problem{9}

证明题

证明：若 $n$ 阶方阵 $A$ 可逆，则

$$
\det(A^{-1}) = \frac{1}{\det(A)}
$$

\solution

由矩阵可逆的定义：$AA^{-1} = I$，其中 $I$ 是单位矩阵。

对等式两边同时取行列式：$\det(AA^{-1}) = \det(I) = 1$。

由行列式的乘法性质 $\det(AB) = \det(A)\det(B)$，得 $\det(A)\det(A^{-1}) = 1$。

因 $A$ 可逆，$\det(A) \neq 0$，故 $\det(A^{-1}) = \dfrac{1}{\det(A)}$。$\blacksquare$

\textbf{Answer:}
\[\boxed{\det(A^{-1}) = \frac{1}{\det(A)}}\]

\vspace{1em}
\hrule
\vspace{1em}


\vspace{2em}
\noindent\textit{Auto-generated: 9/9 problems solved by SymPy.}
\end{document}
